% ============================================================
% Chapter 1 — Introduction to Web Application Security
% ============================================================
\chapter{Introduction to Web Application Security}
\label{ch:intro}

\noteblock{This chapter is designed for absolute beginners. No prior security knowledge is required. Read through carefully before moving to any hands-on chapters.}

% ─────────────────────────────────────────────────────────
\section{What is Web Application Security?}

A \textbf{web application} is any software that runs in a browser — online banking portals, social media platforms, e-commerce stores, admin dashboards, and APIs that power mobile apps. These applications handle highly sensitive data: passwords, payment details, personal information, and confidential business records.

\textbf{Web application security} is the discipline of identifying, understanding, and fixing weaknesses in web applications that could allow an attacker to:

\begin{itemize}[leftmargin=*]
  \item Access data they are not authorized to see
  \item Modify or delete data belonging to others
  \item Impersonate other users (including administrators)
  \item Take over the underlying server
  \item Disrupt the application's availability
\end{itemize}

\noteblock{Think of a web application as a building. Web application security is the process of checking every door, window, lock, and alarm system — before a burglar does.}

% ─────────────────────────────────────────────────────────
\section{Why Organizations Perform Penetration Testing}

Organizations spend significant money building secure applications, yet breaches happen constantly. The reason is simple: \textit{developers are not attackers}. Penetration testing (pen testing) closes this gap by employing skilled professionals who think and act like real attackers — but with written authorization.

\begin{table}[H]
\centering
\caption{Business reasons for penetration testing}
\begin{tabularx}{\textwidth}{>{\bfseries}lX}
\toprule
Reason & Explanation \\
\midrule
Risk Reduction & Find vulnerabilities before real attackers exploit them \\
Compliance & PCI-DSS, ISO 27001, SOC 2 require regular pen tests \\
Customer Trust & Demonstrates commitment to protecting customer data \\
Incident Prevention & A single breach can cost millions and destroy reputation \\
Developer Feedback & Gives development teams real attack scenarios to learn from \\
\bottomrule
\end{tabularx}
\end{table}

% ─────────────────────────────────────────────────────────
\section{Vulnerability Assessment vs Penetration Testing}

These two terms are often confused. They are related but fundamentally different.

\begin{table}[H]
\centering
\caption{Vulnerability Assessment vs Penetration Testing}
\begin{tabularx}{\textwidth}{>{\bfseries}lXX}
\toprule
Attribute & Vulnerability Assessment & Penetration Testing \\
\midrule
Goal & Identify vulnerabilities & Exploit vulnerabilities \\
Depth & Broad scan & Deep, targeted \\
Automation & Mostly automated tools & Manual + automated \\
Risk & Lower (non-exploiting) & Higher (active exploitation) \\
Output & List of vulnerabilities & Proof of exploitation \\
Time & Hours & Days to weeks \\
Use case & Regular scanning & Pre-launch, compliance \\
\bottomrule
\end{tabularx}
\end{table}

\noteblock{A \textbf{vulnerability assessment} tells you \textit{what might be wrong}. A \textbf{penetration test} proves \textit{what can actually be exploited} and shows the real business impact.}

% ─────────────────────────────────────────────────────────
\section{Authorized vs Unauthorized Testing}

\dangerblock{
\textbf{Testing any system without explicit written authorization is illegal.}\\[4pt]
Most countries — including the UK (Computer Misuse Act), USA (CFAA), and EU member states — impose criminal penalties including imprisonment for unauthorized computer access, regardless of intent.\\[6pt]
This guide covers both authorized lab practice and the techniques attackers use in the wild. Understanding unauthorized attack techniques is essential for defenders. However, you must \textbf{only practice on systems you own or have written permission to test}.
}

\subsection{Authorized Testing}

Authorized testing (legitimate penetration testing) always involves:

\begin{enumerate}[leftmargin=*]
  \item \textbf{Rules of Engagement} — written document defining scope, methods, and time window
  \item \textbf{Scope Definition} — exact list of IP addresses, domains, and endpoints in scope
  \item \textbf{Emergency Contacts} — who to call if testing causes unintended impact
  \item \textbf{Safe Harbors} — legal protection clauses for the tester
  \item \textbf{Deliverable Agreement} — what report format and findings will be provided
\end{enumerate}

\subsection{Unauthorized Testing (What NOT to Do)}

Unauthorized testing includes:

\begin{itemize}[leftmargin=*]
  \item Running scanners against websites you do not own
  \item Testing a production system beyond the agreed scope
  \item Using discovered credentials to access other systems
  \item Accessing, copying, or exfiltrating real data — even if accidentally discovered
\end{itemize}

\warnblock{Even if you are ``just testing'' or ``trying to help'', unauthorized testing can result in criminal charges. Always get written authorization first.}

% ─────────────────────────────────────────────────────────
\section{Security Testing Methodology}

Professional penetration testers follow a structured methodology. This guide uses the following phases throughout all attack modules:

\begin{table}[H]
\centering
\caption{Web Application Penetration Testing Phases}
\begin{tabularx}{\textwidth}{>{\bfseries}clX}
\toprule
Phase & Name & Description \\
\midrule
1 & Reconnaissance & Gather information about the target without attacking \\
2 & Scanning \& Enumeration & Map endpoints, parameters, and technologies \\
3 & Vulnerability Identification & Identify potentially exploitable weaknesses \\
4 & Exploitation & Attempt to exploit identified vulnerabilities \\
5 & Post-Exploitation & Understand impact and access depth \\
6 & Reporting & Document all findings with evidence and remediation \\
\bottomrule
\end{tabularx}
\end{table}

% ─────────────────────────────────────────────────────────
\section{OWASP Top 10 Overview}
\label{sec:owasp}

The \textbf{OWASP Top 10} is the most widely referenced list of critical web application security risks, maintained by the Open Web Application Security Project. All attack modules in this guide map to OWASP categories.

\begin{table}[H]
\centering
\caption{OWASP Top 10 (2021)}
\begin{tabularx}{\textwidth}{>{\bfseries}clX}
\toprule
ID & Category & Summary \\
\midrule
A01 & Broken Access Control & Users can act outside their intended permissions \\
A02 & Cryptographic Failures & Weak or missing encryption exposes sensitive data \\
A03 & Injection & Untrusted data sent to an interpreter as a command \\
A04 & Insecure Design & Flaws in the architectural design of the application \\
A05 & Security Misconfiguration & Insecure default configurations, exposed debug info \\
A06 & Vulnerable Components & Using libraries/frameworks with known vulnerabilities \\
A07 & Authentication Failures & Broken login, session management, or MFA issues \\
A08 & Software Integrity Failures & Unverified software updates or CI/CD pipeline attacks \\
A09 & Logging Failures & Insufficient logging prevents detection of attacks \\
A10 & Server-Side Request Forgery & Server fetches attacker-controlled URLs \\
\bottomrule
\end{tabularx}
\end{table}

\subsection{A01 — Broken Access Control}
The most common OWASP finding. Users perform actions or access data beyond their authorized permissions. Examples: accessing \texttt{/admin} without login, reading another user's files by changing an ID parameter.

\subsection{A02 — Cryptographic Failures}
Sensitive data transmitted or stored without proper encryption. Examples: passwords stored as plain MD5, credit card numbers sent over HTTP, weak JWT secrets.

\subsection{A03 — Injection}
The application passes untrusted user input to an interpreter (SQL, OS, LDAP). The interpreter executes the input as a command. SQL Injection is the most well-known example.

\subsection{A04 — Insecure Design}
Security flaws baked into the architecture. No amount of patching fixes a fundamentally broken design. Example: a password reset flow that does not verify ownership of the email address.

\subsection{A05 — Security Misconfiguration}
Default credentials, verbose error messages, exposed debug endpoints, open cloud storage buckets, missing security headers. This is what Level 01 of this course focuses on.

\subsection{A06 — Vulnerable Components}
Using outdated libraries, frameworks, or OS packages with known CVEs. A single vulnerable npm package in a Node.js application can expose the entire server.

\subsection{A07 — Authentication Failures}
Weak passwords permitted, no brute-force protection, session tokens not invalidated on logout, predictable session IDs. Level 02 covers these in depth.

\subsection{A08 — Software Integrity Failures}
Applications that auto-update without verifying signatures, or CI/CD pipelines that can be compromised to inject malicious code into production builds.

\subsection{A09 — Logging \& Monitoring Failures}
Without detailed logs, attackers operate invisibly. Failed login storms, IDOR probes, and SQL injection attempts should all trigger alerts.

\subsection{A10 — Server-Side Request Forgery (SSRF)}
The server makes HTTP requests to attacker-controlled URLs. Attackers use this to reach internal services, AWS metadata endpoints, and internal databases not exposed to the internet.

\exerciseblock{
\textbf{Exercise 1.1 — OWASP Mapping}\\[4pt]
For each of the following scenarios, identify the OWASP A0X:2021 category:
\begin{enumerate}
  \item A URL parameter \texttt{?user\_id=123} can be changed to \texttt{?user\_id=124} to view another user's profile
  \item An application returns a full SQL error including the query text
  \item The \texttt{/admin} panel is accessible without logging in
  \item An old version of \texttt{log4j} is used in the application backend
  \item The application has no rate limiting on the login endpoint
\end{enumerate}
\textit{Answers: IDOR/A01, A05, A01, A06, A07}
}
